\documentclass[conference]{IEEEtran}

\usepackage[utf8]{inputenc}
\usepackage[T1]{fontenc}
\usepackage{cite}
\usepackage{amsmath,amssymb}
\usepackage{booktabs}
\usepackage{array}
\usepackage{xcolor}
\usepackage{listings}
\usepackage{url}

\lstset{
  basicstyle=\ttfamily\footnotesize,
  breaklines=true,
  frame=single,
  columns=fullflexible,
  keepspaces=true,
  showstringspaces=false
}

\title{Milestone 1: Technical Understanding and Research Framing\\Loop Unrolling and Controller Delay in High-Level Synthesis}

\author{\IEEEauthorblockN{Christian Percival}
\IEEEauthorblockA{Electronic Engineering\\
Hochschule Hamm--Lippstadt\\
Lippstadt, Germany\\
Email: percival.christian@gmail.com}}

\begin{document}
\maketitle

\begin{abstract}
This milestone documents the technical understanding and research framing for the seminar topic based on the assigned reference paper by Kurra, Singh, and Panda. The paper studies the impact of loop unrolling on controller delay in high-level synthesis (HLS). The main technical issue is that loop unrolling can reduce latency by exposing parallelism, but it can also increase the complexity and delay of the generated finite-state-machine (FSM) controller. This milestone identifies the problem statement, core technical idea, relevant formalism, assumptions, limitations, resource aspects, implementation direction, and open questions for feedback.
\end{abstract}

\begin{IEEEkeywords}
Milestone 1, high-level synthesis, loop unrolling, controller delay, FSM, timing closure, embedded systems.
\end{IEEEkeywords}

\section{Purpose of Milestone 1}
The purpose of this milestone is not to produce a polished final paper, but to show technical understanding of the assigned reference paper. The milestone focuses on the technical problem, the main contribution, the assumptions, and the research framing. The final paper will later expand this into related work, critical evaluation, and implementation results.

\section{Assigned Reference}
The assigned reference is:

\begin{quote}
S. Kurra, N. K. Singh, and P. R. Panda, ``The Impact of Loop Unrolling on Controller Delay in High Level Synthesis,'' DATE, 2007.
\end{quote}

The paper studies how different loop unroll factors affect controller delay in HLS-generated hardware. It proposes an estimation-based method to choose suitable unroll factors without exhaustively synthesizing all possible versions of the design~\cite{kurra2007impact}.

\section{Problem Statement}
Loop unrolling is a well-known optimization technique. It duplicates the loop body so that multiple iterations can be executed or scheduled together. In software, unrolling is usually used to reduce loop overhead and expose instruction-level parallelism. In HLS, the situation is more complex because the transformed loop is not only compiled into instructions, but synthesized into hardware.

The technical problem is:

\begin{quote}
How can an HLS tool select a loop unroll factor that improves latency without causing excessive FSM/controller delay and timing violations?
\end{quote}

This is important because HLS output usually contains both a datapath and a controller. The datapath contains functional units such as adders, multipliers, registers, memories, and multiplexers. The controller is often an FSM that generates control signals for the datapath. If loop unrolling increases the size and complexity of the FSM too much, the controller can become part of the critical path and limit the clock frequency.

\section{Motivation}
A designer might assume that a larger unroll factor always improves performance. This is not necessarily true in HLS. Unrolling can reduce the number of loop iterations, but it also increases hardware complexity.

Important motivation points:
\begin{itemize}
    \item Loop unrolling can reduce latency by exposing parallelism.
    \item Larger loop bodies can require more operations to be scheduled.
    \item More operations can increase the number of control states and control signals.
    \item More resource sharing can increase multiplexer sizes.
    \item The final design may fail timing if controller delay is ignored.
\end{itemize}

Therefore, loop unrolling should be treated as a design-space exploration problem, not as a simple ``more is better'' optimization.

\section{Core Technical Idea}
The core idea of the reference paper is to estimate the effect of loop unrolling on FSM/controller delay before performing expensive detailed synthesis. The paper proposes selecting unroll factors based on both latency improvement and additional controller delay.

The general optimization goal is:

\begin{quote}
Minimize latency while keeping FSM delay below a given delay budget.
\end{quote}

The paper uses a heuristic similar to a knapsack problem. Each unrolling decision has:
\begin{itemize}
    \item a benefit: reduction in latency;
    \item a cost: increase in FSM/controller delay.
\end{itemize}

A simplified priority function is:

\begin{equation}
Pr(l_i) = \frac{\text{Reduction in latency}}{\text{Additional FSM delay}}.
\label{eq:priority}
\end{equation}

This gives high priority to loops where unrolling gives a large performance benefit with a relatively small controller-delay penalty.

\section{Relevant Formalism and Model}
The reference paper uses high-level information to estimate controller delay. A simplified version of the delay model is:

\begin{equation}
Delay = A \cdot \log(\#operations) \cdot (\#state\ bits),
\label{eq:delay}
\end{equation}

where $A$ is a technology-dependent parameter. It depends on resource constraints, function-unit control bits, and variables. The important idea is that controller delay increases with the number of operations and with the number of bits needed to represent FSM states.

The paper also considers how unrolling changes operation count. If a loop body contains $LOps(l_i)$ operations and is unrolled by factor $u_i$, the number of duplicated loop-body operations increases approximately by:

\begin{equation}
(u_i - 1) \cdot LOps(l_i).
\label{eq:opsincrease}
\end{equation}

If the unroll factor does not evenly divide the loop bound, a trailer loop may be needed. This can reduce the expected latency benefit and add extra operations.

\section{Example Kernel for Understanding}
A simple dot-product loop is a suitable example:

\begin{lstlisting}[language=C,caption={Original loop.}]
#define N 32

int dot_product(int x[N], int z[N]) {
    int q = 0;
    for (int i = 0; i < N; i++) {
        q += x[i] * z[i];
    }
    return q;
}
\end{lstlisting}

A partially unrolled version with factor two is:

\begin{lstlisting}[language=C,caption={Partially unrolled loop.}]
int dot_product(int x[N], int z[N]) {
    int q = 0;
    for (int i = 0; i < N; i += 2) {
        q += x[i] * z[i];
        q += x[i+1] * z[i+1];
    }
    return q;
}
\end{lstlisting}

This version has fewer loop iterations, but the loop body is larger. In an HLS design, this can increase resource demand and controller complexity.

\section{Assumptions and Scope Limitations}
The reference paper makes several assumptions and has scope limitations that are important for the seminar framing.

\subsection{Assumptions}
\begin{itemize}
    \item The input is a behavioral or C-style program that can be processed by an HLS flow.
    \item Loops have analyzable structure and known or estimable iteration counts.
    \item The FSM delay budget is provided as an input to the algorithm.
    \item The delay-estimation constants depend on the target technology and synthesis flow.
    \item The model focuses on controller delay rather than full system-level performance.
\end{itemize}

\subsection{Limitations}
\begin{itemize}
    \item Exact numerical results from the paper are not directly transferable to every FPGA or modern HLS tool.
    \item Modern HLS tools may apply additional optimizations that interact with unrolling.
    \item FPGA timing may be strongly affected by routing, memory architecture, and placement.
    \item Controller delay may not be directly reported by all student-accessible HLS tools.
    \item A low-budget student implementation can demonstrate the trend, but may not reproduce the original ASIC-based delay model exactly.
\end{itemize}

\section{Runtime, Memory, Scalability, and Resource Aspects}
Loop unrolling affects several engineering aspects.

\subsection{Runtime and Latency}
Unrolling can reduce the number of loop iterations and expose parallelism. However, latency does not always decrease monotonically. If the unroll factor does not evenly divide the loop bound, a trailer loop may reduce the benefit.

\subsection{Hardware Resources}
Unrolling may increase the number of adders, multipliers, memory ports, registers, and multiplexers required by the generated hardware. If resources are limited, the HLS scheduler may still serialize operations, reducing the benefit of unrolling.

\subsection{Controller Complexity}
The FSM may require more states, more state bits, and more complex next-state logic. This is the main focus of the reference paper.

\subsection{Scalability}
Exhaustively testing all unroll factors for all loops becomes expensive when a program contains many loops. Therefore, the paper proposes pruning the search space and using estimation instead of full synthesis for every candidate.

\section{Possible Student Implementation Direction}
The implementation will be designed as a small reproducible experiment. The aim is to compare unroll factors for a simple loop kernel and measure or estimate their impact.

\subsection{Option A: HLS-Based Experiment}
Use Vitis HLS, Intel HLS, or another available HLS tool. Synthesize the same loop kernel with unroll factors:

\begin{equation}
U = \{1, 2, 4, 8, 16, 32\}.
\end{equation}

Record:
\begin{itemize}
    \item latency in cycles;
    \item estimated clock period;
    \item timing slack;
    \item LUT/ALM usage;
    \item register usage;
    \item DSP usage;
    \item BRAM or memory usage;
    \item synthesis time.
\end{itemize}

\subsection{Option B: Python Estimation Model}
If HLS tools are not available, implement a simplified model based on the paper's equations. The Python model can estimate latency and controller delay trends for each unroll factor.

\begin{lstlisting}[language=Python,caption={Simplified estimation idea.}]
for u in [1, 2, 4, 8, 16, 32]:
    iterations = ceil(N / u)
    operations = base_ops + (u - 1) * loop_body_ops
    state_bits = ceil(log2(estimated_states))
    controller_delay = A * log2(operations) * state_bits
\end{lstlisting}

The Python model is not a replacement for synthesis, but it can demonstrate why latency and controller delay move in opposite directions.

\section{Expected Output for Later Milestones}
Milestone 1 provides the technical foundation. Later milestones can expand as follows:

\begin{itemize}
    \item Milestone 2: Compare loop unrolling with pipelining, resource sharing, memory partitioning, and design-space exploration.
    \item Milestone 3: Critically evaluate where the paper's approach is useful, where it may fail, and how modern HLS tools change the situation.
    \item Milestone 4: Final paper, slides, implementation results, and discussion.
\end{itemize}

\section{Open Questions for Feedback}
The following questions should be clarified during feedback:

\begin{enumerate}
    \item Should the implementation focus on a real HLS tool, or is a Python estimation model acceptable as the main demonstration?
    \item If using an HLS tool, which tool is preferred by the course or university environment?
    \item Is FPGA timing slack an acceptable indirect measurement if controller delay is not reported separately?
    \item Should the implementation compare unrolling only, or also include pipelining as a related optimization?
    \item How much mathematical detail from the reference paper should be included in the final seminar paper?
\end{enumerate}

\section{Milestone 1 Checklist Mapping}
\begin{table}[htbp]
\caption{Milestone 1 checklist mapping.}
\label{tab:m1checklist}
\centering
\begin{tabular}{p{0.43\linewidth}p{0.45\linewidth}}
\toprule
\textbf{Checklist item} & \textbf{Current status} \\
\midrule
State the technical problem & Covered in Sections III--IV \\
Explain main contribution & Covered in Section V \\
Identify formalism/model & Covered in Section VI \\
Identify assumptions and limitations & Covered in Section VIII \\
Discuss runtime/resource aspects & Covered in Section IX \\
Document unclear points & Covered in Section XI \\
Explain relevance to embedded systems & Covered through HLS, timing, resource, and FPGA discussion \\
\bottomrule
\end{tabular}
\end{table}

\section{Conclusion}
The assigned paper addresses an important HLS optimization problem: loop unrolling can reduce latency, but it can also increase FSM/controller delay. The main contribution is an estimation-based method for choosing unroll factors while respecting a controller-delay budget. For the seminar, this topic is relevant because it connects compiler-style optimization with embedded hardware constraints such as timing closure, resource usage, and scalability.

\begin{thebibliography}{00}
\bibitem{kurra2007impact}
S. Kurra, N. K. Singh, and P. R. Panda, ``The Impact of Loop Unrolling on Controller Delay in High Level Synthesis,'' in \emph{Design, Automation and Test in Europe Conference and Exhibition (DATE)}, 2007.

\bibitem{sarkar2000optimized}
V. Sarkar, ``Optimized unrolling of nested loops,'' in \emph{Proceedings of the 14th International Conference on Supercomputing}, 2000.

\bibitem{carr1994improving}
S. Carr and K. Kennedy, ``Improving the ratio of memory operations to floating-point operations in loops,'' \emph{ACM Transactions on Programming Languages and Systems}, vol. 16, no. 6, 1994.

\bibitem{gupta2006rapid}
G. R. Gupta, M. Gupta, and P. R. Panda, ``Rapid estimation of control delay from high-level specifications,'' in \emph{Proceedings of the Design Automation Conference (DAC)}, 2006, pp. 455--458.

\bibitem{allan1995software}
V. H. Allan, R. B. Jones, R. M. Lee, and S. J. Allan, ``Software pipelining,'' \emph{ACM Computing Surveys}, vol. 27, no. 3, pp. 367--432, 1995.
\end{thebibliography}

\end{document}
